\documentclass[10pt,a4paper]{amsart}
\usepackage[margin=19mm]{geometry}
\usepackage{amsmath,amssymb,mathtools}
\usepackage[scr=rsfs,cal=euler]{mathalfa}
\usepackage{xfrac}
\usepackage[unicode,hidelinks,bookmarks=false]{hyperref}
\newcommand{\ie}{i.e.\ }
\newcommand{\cf}{cf.\ }
\newcommand{\resp}{respectively\ }
\newcommand{\Subsection}{\subsection}
\providecommand{\nobreakdash}{\nobreak\hbox{-}}
\setlength{\emergencystretch}{2em}
\multlinegap0pt
\usepackage{bm}
\usepackage{bbm}
\usepackage{dsfont}
\usepackage{xspace}
\usepackage{cancel}
\usepackage{mathtools}
\usepackage{amsthm}
\makeatletter
\@ifundefined{theorem}{\newtheorem{theorem}{Theorem}}{}
\makeatother
\title[Density and separation for augmented Zarankiewicz numbers]{Density and separation for augmented Zarankiewicz numbers}
\author{Nikita Lebedev}
\date{Source: arXiv:2609.16555v1 (CC BY 4.0)}
\begin{document}
\maketitle

\noindent\emph{Source and licence.} Density and separation for augmented Zarankiewicz numbers. Version arXiv:2609.16555v1 (CC BY 4.0), distributed under the Creative Commons Attribution 4.0 International licence, as recorded in the arXiv licence metadata for this exact version and shown on the abstract page https://arxiv.org/abs/2609.16555v1. Full rights notice: licenses/arxiv-2609.16555-CC-BY-4.0.txt.\par\smallskip

\noindent\textit{Editorial scope.} Counted unit: the Theorem whose source label is \texttt{thm:phasestability} in the article (source0.tex lines 442--457, source label \texttt{thm:phasestability}), delivered together with the article's own complete original proof of it (source0.tex lines 458--646, one \texttt{proof} environment of 189 lines). No mathematics is written, repaired or re-derived: statement and proof are the article's own text line for line. Required context retained uncounted: eq:openpivot (lines 294--297); eq:scalar-slack (lines 299--302); eq:degree-moment (lines 314--316). Every definition, display and label of those lines is retained unchanged. Omitted from this package and not redistributed: 1--293, 298--298, 303--313, 317--441, 647--3036. Nothing inside a retained excerpt is omitted or altered: every retained body is delivered exactly as the article's lines are written, so no omission receipt of any kind accompanies this package. Citations: the retained excerpts carry no citation command at all, so no bibliography entry of the article is transcribed and no reference is redistributed. Reference adaptation: none was needed and none was made; every reference inside the delivered unit resolves inside it, so no unresolved reference or citation is produced. Printed slips kept literal: no printed word, symbol, hypothesis or number of the counted statement or of its proof was altered. Layout and class: the article's own class is \texttt{article} with packages including fontenc, geometry, amsmath, amssymb, amsthm, microtype, array, needspace, tikz, hyperref; the wrapper is the lane's pinned \texttt{amsart} editorial wrapper at 10pt on A4, which restates the theorem-like environments the retained text needs and adds the article's own macro definitions that the retained text needs (none), with extra wrapper packages (bm, bbm, dsfont, xspace, cancel, mathtools). The delivered numbering is the unit's own. Assets: no retained span contains a figure environment, a graphics-inclusion command, a PDF-page inclusion, a table or a TikZ picture; no image file is copied into this package, the package declares no asset dependency and no image file is redistributed with it. Licence: version arXiv:2609.16555v1 of this article is distributed under the Creative Commons Attribution 4.0 International licence, as recorded in the arXiv licence metadata for this exact version and shown on the abstract page https://arxiv.org/abs/2609.16555v1. A full rights notice is delivered at licenses/arxiv-2609.16555-CC-BY-4.0.txt. Characters: every retained excerpt is pure ASCII, so the delivered engine, XeLaTeX with its default Latin Modern text font, reports no missing character.
\begin{equation}\label{eq:openpivot}
 \sum_{\substack{s\ne r\\t_{rs}>0}}
 t_{rs}\beta\bigl((t_{rs}-1)/d_r\bigr)\le E/2.
\end{equation}

\begin{equation}\label{eq:scalar-slack}
 g(u)=2u\beta(u)-3u+1
 =\frac{(1-u)(2u-1)^2}{u^2+(1-u)^2}\ge0.
\end{equation}

\begin{equation}\label{eq:degree-moment}
 3\sum_c e_c^2\le(2m+2)E+3m(m-1),
\end{equation}

\begin{theorem}[Linear stability]\label{thm:phasestability}
There are absolute constants $C,\delta_0>0$ with the following property.
Let all $E$ occupied cells of an $m\times n$ binary matrix be partitioned
into (C3)-pairs. If $\epsilon\ge0$, $E/(mn)\ge2/3-\epsilon$, and
$\delta=\epsilon+1/m+1/n\le\delta_0$, there are row and column partitions
into three parts, with $a_i=|R_i|/m$ and $b_i=|C_i|/n$, such that
\[
 \operatorname{dist}_{\rm edit}
 \left(A,\,\overline{\bigcup_i R_i\times C_i}\right)\le C\delta mn,
 \qquad
 \sum_i(a_i-1/3)^2+\sum_i(b_i-1/3)^2\le C\delta.
\]
In particular, if $m,n\to\infty$ and $E/(mn)\to2/3$, the matrix is
within $o(mn)$ cells of this template with all part sizes asymptotic
to one third of their respective dimensions.
\end{theorem}

\begin{proof}
We first show that most rows have degree close to $2n/3$ and that their
normalized codegrees are close to either $1/2$ or $1$. The root $1$
will identify rows with nearly equal supports; the root $1/2$ will
separate three such classes. The last step turns their approximate
column classes into a partition without losing the linear edit bound.

\smallskip\noindent\emph{Degree concentration and codegree slack.}
Put $u_r=d_r/n$, $v_c=e_c/m$, $\alpha=E/(mn)$, and write $t_{rs}$
for row codegrees. The degree moment and its transpose give
\begin{equation}\label{eq:stabilitymoments}
 \frac1m\sum_r(u_r-2/3)^2
 \le\frac{2\epsilon}{3}+\frac1m+\frac{2}{3n}\le\delta,
 \qquad
 \frac1n\sum_c(v_c-2/3)^2\le\delta.
\end{equation}
Indeed, the transpose of \eqref{eq:degree-moment} gives
$m^{-1}\sum_r u_r^2\le(2/3+2/(3n))\alpha+(n-1)/(mn)$.
Subtracting $4\alpha/3$ and adding $4/9$, then using
$\alpha\ge2/3-\epsilon$ and $\alpha\le1$, proves the first estimate;
the other follows in the same way.

Use $g(q)=2q\beta(q)-3q+1$ from \eqref{eq:scalar-slack}.
For $d_r>0$ put $q'_{rs}=(t_{rs}-1)_+/d_r$. Retaining the slack in
\eqref{eq:openpivot} gives
\[
 G':=\sum_{r:d_r>0}\sum_{s\ne r}d_rg(q'_{rs})
 \le(2m+2)E-3\sum_c e_c^2+3m(m-1)
 \le3\delta m^2n.
\]
Indeed $2t\beta((t-1)_+/d)\ge3t-3-d+dg((t-1)_+/d)$;
at $t=0$ the right side is $-3$, so zero codegrees are included.
Cauchy bounds the preceding right side, divided by $m^2n$, by
$2\alpha-3\alpha^2+2\alpha/m+3(m-1)/(mn)$.
Taking $\delta_0\le1/6$ ensures $\alpha\ge1/2$, where
$2\alpha-3\alpha^2$ is decreasing. It is therefore at most
$2\epsilon-3\epsilon^2$, giving the bound $3\delta$.
Since $0\le\beta\le1$ and
$\beta'(q)=2q(1-q)/(q^2+(1-q)^2)^2\le2$, we have $|g'|\le9<20$.
Removing the unit shift therefore costs at most $20m^2$, and the
unshifted diagonal has $g(1)=0$.
Define the symmetric slack
\[
 S_{rs}=u_rg(t_{rs}/d_r)+u_sg(t_{rs}/d_s),
 \qquad \frac1{m^2}\sum_{r,s}S_{rs}\le46\delta,
\]
with zero-degree summands interpreted as zero.
Splitting at $3/4$ in the formula for $g$ shows
\begin{equation}\label{eq:stabilityscalar}
 g(q)\ge\operatorname{dist}(q,\{1/2,1\})^2,
 \qquad g(q)\ge(1-q)/4\quad(3/4\le q\le1).
\end{equation}

\smallskip\noindent\emph{Near classes and their sizes.}
Fix a large absolute $A$. Call a row a representative if
\[
 |u_r-2/3|\le\sqrt{A\delta},\qquad
 \frac1m\sum_sS_{rs}\le A\delta.
\]
Markov's inequality applied to the two nonnegative averages above bounds
the fractions failing these tests by $1/A$ and $46/A$.
Choose $A$ so their sum $47/A$ is less than $1/100$.
Shrink $\delta_0$ so representatives have $u_r\ge3/5$.
For each representative,
\begin{equation}\label{eq:representativeslack}
 \frac1m\sum_s\operatorname{dist}(t_{rs}/d_r,\{1/2,1\})^2
 \le2A\delta.
\end{equation}
Put $\tau=1/20$ and
$L(r)=\{s:|N_r\mathbin\triangle N_s|\le\tau n\}$.
For $s\in L(r)$, each of $d_r-t_{rs}$ and $d_s-t_{rs}$ is at most
$\tau n$, while $d_s\ge d_r-\tau n\ge11n/20$.
Thus both codegree ratios exceed $3/4$, so
\eqref{eq:stabilityscalar} gives
\begin{equation}\label{eq:lineareditnear}
 |N_r\mathbin\triangle N_s|/n\le4S_{rs},\qquad
 \sum_{s\in L(r)}|N_r\mathbin\triangle N_s|\le4A\delta mn.
\end{equation}
Here $|N_r\mathbin\triangle N_s|=d_r+d_s-2t_{rs}$: applying
the second scalar bound to the two summands in $S_{rs}$ controls each
of $d_r-t_{rs}$ and $d_s-t_{rs}$. This linear estimate is what will
keep the final edit error at order $\delta$, rather than $\sqrt\delta$.

Fix $\zeta=1/1000$. For each representative $r$, Markov applied to
\eqref{eq:stabilitymoments} and \eqref{eq:representativeslack} shows
that, apart from $O(\delta)m$ rows $s$, both
$|u_s-2/3|\le\zeta$ and
$\operatorname{dist}(t_{rs}/d_r,\{1/2,1\})\le\zeta$ hold.
Take $\delta_0$ small enough that the representative's degree is also
within $\zeta$ of $2/3$. Write $q=t_{rs}/d_r$. If $q\ge1-\zeta$, then
$|N_r\mathbin\triangle N_s|/n=u_s-u_r+2u_r(1-q)\le4\zeta<\tau$.
If $|q-1/2|\le\zeta$, the same distance is at least
$u_s-2\zeta u_r\ge2/3-3\zeta>\tau$.
Thus the root-one alternative places $s$ in $L(r)$, and the
root-half alternative places it outside. For the row zero sets
$Z_r=[n]\setminus N_r$ and $Z_s=[n]\setminus N_s$, the latter gives
\[
 |Z_r\cap Z_s|/n\le1-u_s-(1/2-\zeta)u_r
       \le13\zeta/6-\zeta^2<3\zeta.
\]
For $\theta=|L(r)|/m$, equation~\eqref{eq:lineareditnear} gives
\[
 \frac1m\sum_{s\in L(r)}\left|\frac{t_{rs}}n-u_r\right|=O(\delta).
\]
Outside $L(r)$, all but $O(\delta)m$ rows take the half-root alternative.
On those rows $|q-1/2|=\operatorname{dist}(q,\{1/2,1\})$; Cauchy and
\eqref{eq:representativeslack} bound its average by $\sqrt{2A\delta}$.
The exceptional rows contribute $O(\delta)$ because all normalized
codegrees lie in $[0,1]$. Consequently
\[
 \frac{\sum_s t_{rs}}{mn}
 =\frac{u_r(1+\theta)}2+O(\sqrt\delta).
\]
On the other hand, interchanging the codegree sum gives
\[
 \frac{\sum_s t_{rs}}{mn}
 =\frac1n\sum_{c\in N_r}v_c
 =\frac{2u_r}{3}+O(\sqrt\delta),
\]
since Cauchy bounds the last error by
$\sqrt{u_r\,n^{-1}\sum_c(v_c-2/3)^2}\le\sqrt\delta$.
Comparing the two expressions and using $u_r\ge3/5$ proves
$|L(r)|/m=1/3+O(\sqrt\delta)$ for every representative.

\smallskip\noindent\emph{Three prototypes cover almost all rows.}
Choose three representatives successively outside the previous near
classes. Also require, from each previously chosen representative's
perspective, codegree distance to $\{1/2,1\}$ at most $B\sqrt\delta$.
Each extra test excludes at most $2A/B^2$ of the rows. Choose $B$ after
$A$, large enough that $4A/B^2<1/100$. At the third choice the two
near classes exclude at most $2/3+O(\sqrt\delta)$ of the rows, and
nonrepresentatives and the two extra tests exclude less than $2/100$.
Thus all three choices are possible after shrinking $\delta_0$.
Their mutual codegrees must be near the half root, since the
one root would put them in the same near class. Take $\delta_0$ small
enough that $\max\{B,\sqrt A\}\sqrt\delta\le\zeta$. For two representatives $r_i,r_j$ the
half root then gives
$|N_{r_i}\mathbin\triangle N_{r_j}|\ge d_{r_j}-2\zeta d_{r_i}\ge(2/3-3\zeta)n>2\tau n$,
so their near classes are disjoint by the Hamming triangle inequality.
Put $Z_i=Z_{r_i}$. The stronger, $B\sqrt\delta$ tests give
$t_{r_i r_j}/n=1/3+O(\sqrt\delta)$, and hence
$|Z_i\cap Z_j|/n=1-u_{r_i}-u_{r_j}+t_{r_i r_j}/n=O(\sqrt\delta)$.
Together with the prototype degrees and inclusion--exclusion, this proves
\begin{equation}\label{eq:prototypezeros}
 |Z_i|=n/3+O(\sqrt\delta)n,\quad
 |Z_i\cap Z_j|=O(\sqrt\delta)n,\quad
 |\bigcup_iZ_i|\ge n-O(\sqrt\delta)n.
\end{equation}

Only $O(\delta)m$ rows lie outside all three near classes.
To see this, discard rows failing either fixed $\zeta$ test for any
prototype, costing $O(\delta)m$ rows. A remaining row outside all near
classes would satisfy $|Z_s\cap Z_i|\le3\zeta n$ for each $i$.
Its intersection with the union of the prototypes is therefore at most
$9\zeta n$, and at most $O(\sqrt\delta)n$ further columns lie outside
that union. By \eqref{eq:prototypezeros} this would give
\[
 (1/3-\zeta)n\le|Z_s|\le9\zeta n+O(\sqrt\delta)n,
\]
which is impossible for small $\delta_0$.
Assign the leftover rows arbitrarily to the three disjoint near classes,
and copy each prototype to its resulting row class, obtaining $A_0$.
The three near classes cost $O(\delta)mn$ edits by
\eqref{eq:lineareditnear}; the leftover rows cost at most $n$ edits each.
Thus
$\operatorname{dist}_{\rm edit}(A,A_0)=O(\delta)mn$ and
$a_i=1/3+O(\sqrt\delta)$.

\smallskip\noindent\emph{Recovering a column partition.}
The sets $Z_i$ need not yet form a partition. A column of $A_0$ has normalized degree
$v_0(c)=1-\sum_{i:c\in Z_i}a_i$.
If it belongs to zero, two or three prototype zero sets, this degree is
within $O(\sqrt\delta)$ of $1$, $1/3$ or $0$, respectively, and hence
at least $1/5$ from $2/3$ for small $\delta_0$.
Only $O(\delta)n$ original columns are at
least $1/10$ from $2/3$, by \eqref{eq:stabilitymoments}.
The other columns with an invalid zero-set membership count have
$|v_0(c)-v_c|\ge1/10$; there are only $O(\delta)n$ of them because
$\sum_c|v_0(c)-v_c|$ is bounded by the edit count divided by $m$.
Assign these exceptional columns to single zero-set labels. This produces
a partition $C_i$ by changing $O(\delta)n$ columns and costs another
$O(\delta)mn$ edits. Each $C_i$ differs from $Z_i$ in $O(\delta)n$ places,
so \eqref{eq:prototypezeros} gives $b_i=1/3+O(\sqrt\delta)$.
The two squared mass bounds follow by summing the six squared
$O(\sqrt\delta)$ deviations. The choices were made in the order
$A$, then $B$, then a sufficiently small $\delta_0$; $\tau$ and $\zeta$
were fixed numerical constants throughout. Thus one absolute choice of
$\delta_0$ suffices for the whole proof.
\end{proof}

\end{document}
